% ============================================================
%  PEMBAHASAN SUPER DETAIL — LATIHAN SOAL & STUDI KASUS
%  BAB 6: FUNGSI KUADRAT
%  Catatan: Soal Akhir Bab TIDAK dibahas di sini (sesuai permintaan).
% ============================================================
% ============================================================
%  PREAMBUL BERSAMA — BUKU PEMBAHASAN LATIHAN SOAL
%  Dipakai oleh MAE-2026-2027-Latihan1.tex ... Latihan9.tex
%  Kompilasi: pdfLaTeX (disarankan Overleaf / MiKTeX)
% ============================================================
\documentclass[11pt,a4paper]{report}

% ---------- PAKET ----------
\usepackage[utf8]{inputenc}
\usepackage[T1]{fontenc}
\usepackage[bahasa]{babel}
\usepackage{amsmath,amssymb,amsthm,mathtools}
\usepackage{geometry}
\usepackage{xcolor}
\usepackage{graphicx}
\usepackage{enumitem}
\usepackage{tikz}
\usetikzlibrary{calc,arrows.meta,angles,quotes,positioning,decorations.pathreplacing,patterns}
\usepackage{pgfplots}
\pgfplotsset{compat=1.18}
\usepackage[most]{tcolorbox}
\usepackage{fancyhdr}
\usepackage{titlesec}
\usepackage{array,booktabs}
\usepackage{multicol}
\usepackage{pifont}
\usepackage[colorlinks=true,linkcolor=biru,urlcolor=biru]{hyperref}

\geometry{a4paper,margin=2.5cm,headheight=15pt}

% ---------- WARNA ----------
\definecolor{biru}{HTML}{1F4E79}
\definecolor{birumuda}{HTML}{D6E4F0}
\definecolor{hijau}{HTML}{2E7D32}
\definecolor{hijaumuda}{HTML}{E3F2E1}
\definecolor{oranye}{HTML}{C75B12}
\definecolor{oranyemuda}{HTML}{FBE8DA}
\definecolor{abu}{HTML}{F2F2F2}
\definecolor{ungu}{HTML}{6A1B9A}
\definecolor{ungumuda}{HTML}{EDE2F3}
\definecolor{merah}{HTML}{C62828}
\definecolor{merahmuda}{HTML}{FCE4E4}
\definecolor{tosca}{HTML}{00838F}
\definecolor{toscamuda}{HTML}{E0F2F1}
\definecolor{emas}{HTML}{8D6E00}
\definecolor{emasmuda}{HTML}{FFF6D6}

% ---------- HEADER / FOOTER ----------
\newcommand{\babjudul}{Pembahasan Latihan}
\pagestyle{fancy}
\fancyhf{}
\renewcommand{\headrulewidth}{0.8pt}
\renewcommand{\footrulewidth}{0pt}
\fancyhead[L]{\small\textcolor{ungu}{\leftmark}}
\fancyhead[R]{\small\textcolor{ungu}{\babjudul}}
\fancyfoot[C]{\thepage}

% ---------- GAYA JUDUL BAB ----------
\titleformat{\chapter}[display]
  {\normalfont\huge\bfseries\color{ungu}}
  {\filright\large PEMBAHASAN}{8pt}{\Huge}
\titlespacing*{\chapter}{0pt}{0pt}{20pt}

% ---------- LINGKUNGAN TEOREMA ----------
\newtheoremstyle{gaya}{6pt}{6pt}{}{}{\bfseries\color{biru}}{.}{.5em}{}
\theoremstyle{gaya}
\newtheorem{definisi}{Definisi}[chapter]
\newtheorem{sifat}{Sifat}[chapter]

% ---------- KOTAK: SOAL (ungu) ----------
\newtcolorbox{soalbox}[1]{
  enhanced, breakable, colback=ungumuda, colframe=ungu, boxrule=0.8pt,
  arc=3pt, left=10pt, right=10pt, top=8pt, bottom=8pt,
  title={\textbf{Soal #1}}, fonttitle=\bfseries\color{white}, coltitle=white,
  attach boxed title to top left={xshift=10pt,yshift=-10pt},
  boxed title style={colback=ungu,arc=2pt}}

% ---------- KOTAK: KONSEP KUNCI (hijau) ----------
\newtcolorbox{ingat}[1]{
  enhanced, breakable, colback=hijaumuda, colframe=hijau, boxrule=0.6pt,
  arc=3pt, left=10pt, right=10pt, top=6pt, bottom=6pt,
  title={\textbf{Konsep Kunci: #1}}, fonttitle=\bfseries\color{white}, coltitle=white,
  attach boxed title to top left={xshift=10pt,yshift=-10pt},
  boxed title style={colback=hijau,arc=2pt}}

% ---------- KOTAK: JEBAKAN (merah) ----------
\newtcolorbox{jebakan}{
  enhanced, breakable, colback=merahmuda, colframe=merah, boxrule=0.6pt,
  arc=3pt, left=10pt, right=10pt, top=6pt, bottom=6pt,
  title={\textbf{\ding{55}\ Jebakan \& Kesalahan yang Sering Terjadi}},
  fonttitle=\bfseries\color{white}, coltitle=white,
  attach boxed title to top left={xshift=10pt,yshift=-10pt},
  boxed title style={colback=merah,arc=2pt}}

% ---------- KOTAK: VARIASI SOAL (oranye) ----------
\newtcolorbox{variasi}{
  enhanced, breakable, colback=oranyemuda, colframe=oranye, boxrule=0.6pt,
  arc=3pt, left=10pt, right=10pt, top=6pt, bottom=6pt,
  title={\textbf{\ding{72}\ Bentuk Modifikasi Soal Ini}},
  fonttitle=\bfseries\color{white}, coltitle=white,
  attach boxed title to top left={xshift=10pt,yshift=-10pt},
  boxed title style={colback=oranye,arc=2pt}}

% ---------- KOTAK: STUDI KASUS (biru) ----------
\newtcolorbox{studikasus}[1]{
  enhanced, breakable, colback=abu, colframe=biru, boxrule=0.8pt,
  arc=3pt, left=10pt, right=10pt, top=8pt, bottom=8pt,
  title={\textbf{Studi Kasus: #1}}, fonttitle=\bfseries\color{white}, coltitle=white,
  attach boxed title to top left={xshift=10pt,yshift=-10pt},
  boxed title style={colback=biru,arc=2pt}}

% ---------- KOTAK: TIPS (tosca) ----------
\newtcolorbox{tips}{
  enhanced, breakable, colback=toscamuda, colframe=tosca, boxrule=0.6pt,
  arc=3pt, left=10pt, right=10pt, top=6pt, bottom=6pt,
  borderline west={3pt}{0pt}{tosca}}

% ---------- PERINTAH BANTU ----------
\newcommand{\jawab}{\par\smallskip\noindent\textit{\textcolor{oranye}{Penyelesaian langkah demi langkah:}}\par\smallskip}
\newcommand{\langkah}[1]{\par\smallskip\noindent\textbf{\textcolor{biru}{Langkah #1.}}\ }
\newcommand{\jadi}{\par\smallskip\noindent$\Rightarrow$\ \textbf{Jadi, }}
% Judul tiap nomor soal
\newcommand{\soalno}[2]{%
  \par\bigskip
  \noindent{\large\bfseries\color{ungu}\rule[-2pt]{4pt}{14pt}\ Soal #1}%
  \ {\normalfont\itshape\color{gray}(#2)}\par\nopagebreak\smallskip}

% ---------- HALAMAN JUDUL OTOMATIS ----------
% #1 = judul topik (mis. "Barisan dan Deret"); #2 = nomor bab (angka)
\newcommand{\buathalamanjudul}[2]{%
  \begin{titlepage}
  \centering
  \vspace*{2cm}
  {\color{ungu}\rule{\textwidth}{2pt}}\\[0.6cm]
  {\Huge\bfseries\color{ungu} PEMBAHASAN\\[0.3cm] LATIHAN SOAL}\\[0.4cm]
  {\Large BAB #2 --- #1}\\[0.2cm]
  {\color{ungu}\rule{\textwidth}{2pt}}\\[1.2cm]

  {\large Matematika SMA --- Fase E (Kelas 10)}\\[0.2cm]
  {\large Kurikulum Merdeka}\\[2cm]

  \begin{tcolorbox}[enhanced,width=0.85\textwidth,colback=ungumuda,colframe=ungu,
    arc=4pt,boxrule=1pt]
  \centering\small
  Buku pembahasan ini mengupas \textbf{setiap soal pada kotak ungu (Latihan Soal)}
  di tiap subbab, \emph{selangkah demi selangkah seperti mengajari dari nol}:
  konsep yang dipakai, jebakan yang sering menjerat, cara soal bisa dimodifikasi,
  dan visualisasi. \emph{Soal Akhir Bab tidak dibahas di sini.}
  \end{tcolorbox}

  \vfill
  {\large Tahun Pelajaran 2026/2027}
  \end{titlepage}
  \tableofcontents
  \thispagestyle{empty}
  \clearpage}

% ============================================================
%  PENJAGA KOMPILASI
%  File ini adalah PREAMBUL BERSAMA (di-\input oleh Latihan1..9),
%  BUKAN untuk dikompilasi sendiri. Jika editor terlanjur mengompilasi
%  file ini secara langsung, tampilkan satu halaman info (bukan error).
%  Saat di-\input oleh berkas Latihan, blok ini dilewati.
% ============================================================
\ifnum\pdfstrcmp{\jobname}{MAE-2026-2027-Latihan-preamble}=0
  \begin{document}
  \thispagestyle{empty}
  \begin{center}
  \vspace*{3cm}
  {\Huge\bfseries\color{ungu} Berkas Preambel Bersama}\\[1cm]
  \begin{tcolorbox}[enhanced,width=0.85\textwidth,colback=ungumuda,
    colframe=ungu,arc=4pt,boxrule=1pt]
  \centering
  File ini \textbf{bukan untuk dikompilasi sendiri.} Ia berisi pengaturan
  bersama (paket, warna, kotak, perintah) yang dipanggil lewat
  \texttt{\textbackslash input} oleh berkas \textbf{MAE-2026-2027-Latihan1.tex}
  sampai \textbf{Latihan9.tex}.

  \medskip
  Silakan buka dan kompilasi salah satu berkas \textbf{Latihan1--Latihan9}
  tersebut untuk menghasilkan PDF pembahasan.
  \end{tcolorbox}
  \end{center}
  \end{document}
\fi

\renewcommand{\babjudul}{Pembahasan Latihan — Bab 6}

\begin{document}
\setcounter{chapter}{6}
\buathalamanjudul{Fungsi Kuadrat}{6}

% ============================================================
\chapter*{Pembahasan Bab 6: Fungsi Kuadrat}
\addcontentsline{toc}{chapter}{Pembahasan Bab 6: Fungsi Kuadrat}
\markboth{PEMBAHASAN BAB 6}{}
% ============================================================

\begin{tips}
\textbf{Cara membaca buku ini.} Tiap soal: \textbf{(1)} soal ditulis ulang,
\textbf{(2)} konsep kunci, \textbf{(3)} langkah demi langkah, \textbf{(4)}
jebakan, \textbf{(5)} variasi. \emph{Aturan emas Bab 6:} sebisa mungkin
\textbf{gambar parabolanya} (arah buka, puncak, titik potong) --- banyak jebakan
tanda lenyap begitu kamu melihat grafiknya.
\end{tips}

% ============================================================
\section*{Studi Kasus: Air Mancur dan Lintasan Proyektil}
\addcontentsline{toc}{section}{Studi Kasus: Air Mancur}
% ============================================================

\begin{studikasus}{Air Mancur dan Lintasan Proyektil}
Ketinggian air (m) terhadap jarak mendatar $x$ (m):
\begin{center}
\begin{tabular}{c|ccccc}
\toprule
$x$ (m) & 0 & 1 & 2 & 3 & 4 \\ \midrule
tinggi (m) & 0 & 3 & 4 & 3 & 0 \\ \bottomrule
\end{tabular}
\end{center}
\textbf{Pertanyaan.} (a) Di mana titik tertinggi? (b) Apakah data simetris?
Terhadap garis apa? (c) Dapatkah kamu menebak rumus tingginya?
\end{studikasus}

\subsection*{Jawaban lengkap studi kasus}

\begin{ingat}{Ciri parabola pada data}
Data yang ``naik lalu turun secara simetris'' adalah tanda fungsi kuadrat. Titik
tertinggi/terendah = \emph{puncak}; garis vertikal yang melewati puncak =
\emph{sumbu simetri}.
\end{ingat}

\textbf{(a) Titik tertinggi.} Nilai terbesar pada tabel adalah $4$ m, tercapai
saat $x=2$. Jadi \textbf{puncak $(2,4)$}.

\textbf{(b) Kesimetrian.} Ya. Perhatikan pasangan yang berjarak sama dari $x=2$:
$x=1$ dan $x=3$ sama-sama $3$; $x=0$ dan $x=4$ sama-sama $0$. Datanya simetris
terhadap garis \textbf{$x=2$} (sumbu simetri, melewati puncak).

\textbf{(c) Menebak rumus.} Gunakan \emph{bentuk puncak} $f(x)=a(x-2)^2+4$
(karena puncak $(2,4)$). Cari $a$ dengan titik lain, mis.\ $(0,0)$:
\[
0=a(0-2)^2+4=4a+4 \Rightarrow a=-1.
\]
Maka $f(x)=-(x-2)^2+4=-x^2+4x$. \emph{Cek} $(1,3)$: $-1+4=3$ \checkmark; $(3,3)$:
$-9+12=3$ \checkmark.
\jadi rumus tinggi air mancur $f(x)=-x^2+4x$, puncak $(2,4)$, simetris terhadap
$x=2$.

% ############################################################
\section*{Subbab 6.1 — Bentuk, Grafik, dan Karakteristik}
\addcontentsline{toc}{section}{Subbab 6.1 — Bentuk \& Grafik}
% ############################################################

\begin{ingat}{Empat fakta parabola $f(x)=ax^2+bx+c$}
\begin{itemize}[leftmargin=*,topsep=1pt]
  \item $a>0$ buka \emph{atas} (punya minimum); $a<0$ buka \emph{bawah} (punya
  maksimum).
  \item Sumbu simetri: $x=-\dfrac{b}{2a}$ (tanda \emph{minus}!).
  \item Puncak: $\left(-\dfrac{b}{2a},\ -\dfrac{D}{4a}\right)$, $D=b^2-4ac$.
  \item Potong sumbu-$y$: $(0,c)$; potong sumbu-$x$: akar $ax^2+bx+c=0$.
  \item Diskriminan: $D>0$ dua akar, $D=0$ satu (kembar), $D<0$ tak ada akar real.
\end{itemize}
\end{ingat}

% ---------------- SOAL 6.1.1 ----------------
\soalno{6.1.1}{jenis akar via diskriminan}
\begin{soalbox}{}
Tentukan jenis akar dari $2x^2+3x+5=0$ tanpa menghitung akarnya.
\end{soalbox}

\jawab
\langkah{1} Kenali $a=2,\ b=3,\ c=5$.
\langkah{2} Hitung diskriminan: $D=b^2-4ac=3^2-4(2)(5)=9-40=-31$.
\langkah{3} Karena $D<0$, tak ada akar real.
\jadi persamaan ini \textbf{tidak punya akar real} (grafiknya tidak memotong
sumbu-$x$).

\begin{jebakan}
\begin{itemize}[leftmargin=*,topsep=2pt]
  \item \textbf{Salah tanda $4ac$.} $4\cdot2\cdot5=40$, dan $D=9-40$, bukan
  $9+40$.
  \item \textbf{Mengira $D<0$ berarti ``tak ada penyelesaian apa pun''.} Grafik
  tetap ada; hanya saja tak menyentuh sumbu-$x$.
\end{itemize}
\end{jebakan}

\begin{variasi}
Menentukan nilai $c$ atau $m$ agar $D>0$/$D=0$/$D<0$ (soal ``syarat akar'').
\end{variasi}

% ---------------- SOAL 6.1.2 ----------------
\soalno{6.1.2}{titik puncak}
\begin{soalbox}{}
Tentukan titik puncak $f(x)=x^2-6x+5$.
\end{soalbox}

\jawab
\langkah{1} $a=1,\ b=-6,\ c=5$. Absis puncak (sumbu simetri):
$x=-\dfrac{b}{2a}=-\dfrac{-6}{2}=3$.
\langkah{2} Ordinat puncak: substitusi $x=3$:
$f(3)=3^2-6(3)+5=9-18+5=-4$.
\jadi titik puncak $(3,-4)$.

\begin{jebakan}
\begin{itemize}[leftmargin=*,topsep=2pt]
  \item \textbf{Salah tanda $-\frac{b}{2a}$.} $b=-6$, jadi $-\frac{-6}{2}=+3$
  (bukan $-3$).
  \item \textbf{Menghitung ordinat pakai $-\frac{D}{4a}$ tapi salah $D$.}
  Substitusi langsung ke $f(x)$ lebih aman.
\end{itemize}
\end{jebakan}

\begin{variasi}
Ubah ke bentuk puncak dengan melengkapkan kuadrat:
$f=(x-3)^2-4$, membaca puncak langsung.
\end{variasi}

% ---------------- SOAL 6.1.3 ----------------
\soalno{6.1.3}{akar dengan pemfaktoran}
\begin{soalbox}{}
Tentukan akar $x^2-5x+6=0$.
\end{soalbox}

\jawab
\langkah{1} Cari dua bilangan yang \emph{dikali} $=6$ (yakni $c$) dan
\emph{dijumlah} $=-5$ (yakni $b$): keduanya $-2$ dan $-3$.
\langkah{2} Faktorkan: $(x-2)(x-3)=0$.
\langkah{3} Akar: $x=2$ atau $x=3$.
\jadi $x=2$ atau $x=3$.

\begin{jebakan}
\begin{itemize}[leftmargin=*,topsep=2pt]
  \item \textbf{Salah tanda faktor.} $(x-2)(x-3)$ memberi $+6$ dan $-5$; jangan
  $(x+2)(x+3)$ (itu $+6$ dan $+5$).
\end{itemize}
\end{jebakan}

\begin{variasi}
Jika tak bisa difaktorkan rapi, pakai rumus $abc$; hubungkan dengan Vieta
($x_1+x_2=5,\ x_1x_2=6$).
\end{variasi}

% ---------------- SOAL 6.1.4 ----------------
\soalno{6.1.4}{menggambar parabola buka bawah}
\begin{soalbox}{}
Gambarkan grafik $f(x)=-x^2+2x+3$ lengkap dengan titik puncak.
\end{soalbox}

\jawab
\langkah{1} $a=-1<0\Rightarrow$ buka ke \emph{bawah} (punya maksimum).
\langkah{2} Puncak: $x=-\dfrac{b}{2a}=-\dfrac{2}{2(-1)}=1$; $f(1)=-1+2+3=4$.
Puncak $(1,4)$.
\langkah{3} Titik potong sumbu-$x$: $-x^2+2x+3=0\Rightarrow x^2-2x-3=0\Rightarrow
(x-3)(x+1)=0\Rightarrow x=3,\ -1$.
\langkah{4} Titik potong sumbu-$y$: $(0,3)$.

\begin{center}
\begin{tikzpicture}
\begin{axis}[axis lines=middle,width=8.5cm,height=5.5cm,xlabel=$x$,ylabel=$y$,
  xmin=-2.5,xmax=4.5,ymin=-1.5,ymax=5,samples=60]
\addplot[very thick,biru,domain=-1.6:3.6]{-x^2+2*x+3};
\addplot[only marks,mark=*,oranye] coordinates {(1,4)};
\node[oranye,font=\scriptsize,above] at (axis cs:1,4){puncak $(1,4)$};
\addplot[only marks,mark=*,hijau] coordinates {(3,0)(-1,0)};
\end{axis}
\end{tikzpicture}
\end{center}
\jadi parabola buka bawah, puncak $(1,4)$, memotong sumbu-$x$ di $-1$ dan $3$,
sumbu-$y$ di $3$.

\begin{jebakan}
\begin{itemize}[leftmargin=*,topsep=2pt]
  \item \textbf{Menggambar buka atas.} $a=-1<0$ harus buka ke \emph{bawah}.
  \item \textbf{Lupa mengalikan $-1$} saat menyelesaikan $-x^2+2x+3=0$; balik jadi
  $x^2-2x-3=0$ agar mudah difaktorkan.
\end{itemize}
\end{jebakan}

\begin{variasi}
$f(x)=x^2-2x-3$ (buka atas, puncak $(1,-4)$) --- pencerminan.
\end{variasi}

% ---------------- SOAL 6.1.5 ----------------
\soalno{6.1.5}{sumbu simetri}
\begin{soalbox}{}
Tentukan sumbu simetri $f(x)=2x^2-8x+1$.
\end{soalbox}

\jawab
\langkah{1} $a=2,\ b=-8$. Sumbu simetri $x=-\dfrac{b}{2a}=-\dfrac{-8}{2\cdot2}
=\dfrac{8}{4}=2$.
\jadi sumbu simetri $x=2$.

\begin{jebakan}
\begin{itemize}[leftmargin=*,topsep=2pt]
  \item \textbf{Lupa $2a$ di penyebut} ($=4$, bukan $2$).
\end{itemize}
\end{jebakan}

\begin{variasi}
Cari nilai minimum (substitusi $x=2$): $f(2)=8-16+1=-7$.
\end{variasi}

% ---------------- SOAL 6.1.6 ----------------
\soalno{6.1.6}{syarat akar kembar}
\begin{soalbox}{}
Tentukan nilai $c$ agar $x^2+4x+c=0$ punya akar kembar.
\end{soalbox}

\begin{ingat}{Akar kembar $\iff D=0$}
Satu akar (kembar) terjadi tepat saat diskriminan nol. Gunakan $D=b^2-4ac=0$
sebagai persamaan untuk mencari parameter.
\end{ingat}

\jawab
\langkah{1} $a=1,\ b=4$. Syarat kembar: $D=b^2-4ac=0$.
\langkah{2} $16-4(1)c=0\Rightarrow 16=4c\Rightarrow c=4$.
\jadi $c=4$ (persamaannya $x^2+4x+4=(x+2)^2=0$, akar kembar $x=-2$).

\begin{jebakan}
\begin{itemize}[leftmargin=*,topsep=2pt]
  \item \textbf{Memakai $D>0$.} Akar \emph{kembar} butuh $D=0$ (bukan $>0$).
\end{itemize}
\end{jebakan}

\begin{variasi}
Cari $m$ pada $x^2-mx+9=0$ (Soal 6.2.8); atau syarat \emph{dua akar berbeda}
($D>0$).
\end{variasi}

% ---------------- SOAL 6.1.7 ----------------
\soalno{6.1.7}{menyusun fungsi dari akar}
\begin{soalbox}{}
Sebuah fungsi kuadrat memotong sumbu-$x$ di $2$ dan $5$. Tulis salah satu
bentuknya.
\end{soalbox}

\begin{ingat}{Bentuk akar (faktor)}
Jika akarnya $p$ dan $q$, maka $f(x)=a(x-p)(x-q)$ untuk sembarang $a\neq0$.
Karena $a$ bebas, ada tak hingga fungsi --- ``salah satu'' cukup ambil $a=1$.
\end{ingat}

\jawab
\langkah{1} Akar $2$ dan $5$ $\Rightarrow$ faktor $(x-2)$ dan $(x-5)$.
\langkah{2} Ambil $a=1$: $f(x)=(x-2)(x-5)=x^2-7x+10$.
\jadi salah satu bentuknya $f(x)=x^2-7x+10$.

\begin{jebakan}
\begin{itemize}[leftmargin=*,topsep=2pt]
  \item \textbf{Salah tanda faktor.} Akar $2$ berasal dari $(x-2)$, bukan
  $(x+2)$.
\end{itemize}
\end{jebakan}

\begin{variasi}
Jika diberi titik tambahan (mis.\ melalui $(0,20)$), nilai $a$ terkunci; atau
puncak diketahui $\Rightarrow$ pakai bentuk puncak.
\end{variasi}

% ---------------- SOAL 6.1.8 ----------------
\soalno{6.1.8}{nilai minimum}
\begin{soalbox}{}
Tentukan nilai minimum $f(x)=x^2+2x+5$.
\end{soalbox}

\jawab
\langkah{1} $a=1>0\Rightarrow$ punya minimum di puncak. Absis:
$x=-\dfrac{2}{2}=-1$.
\langkah{2} Nilai minimum $=f(-1)=(-1)^2+2(-1)+5=1-2+5=4$.
\jadi nilai minimum $=4$ (tercapai di $x=-1$).

\begin{jebakan}
\begin{itemize}[leftmargin=*,topsep=2pt]
  \item \textbf{Mengira $c=5$ adalah minimum.} $c$ hanya nilai di $x=0$, bukan
  puncak. Minimum di $x=-1$.
\end{itemize}
\end{jebakan}

\begin{variasi}
Melengkapkan kuadrat: $f=(x+1)^2+4$, membaca minimum $4$ langsung.
\end{variasi}

% ---------------- SOAL 6.1.9 ----------------
\soalno{6.1.9}{titik potong sumbu-y}
\begin{soalbox}{}
Tentukan titik potong sumbu-$y$ dari $f(x)=3x^2-x+2$.
\end{soalbox}

\jawab
\langkah{1} Sumbu-$y$ berarti $x=0$. Substitusi: $f(0)=3(0)-0+2=2$.
\jadi titik potong sumbu-$y$ adalah $(0,2)$ (yaitu $c$).

\begin{jebakan}
\begin{itemize}[leftmargin=*,topsep=2pt]
  \item \textbf{Menukar sumbu.} Potong sumbu-$y$: set $x=0$. Potong sumbu-$x$:
  set $y=0$.
\end{itemize}
\end{jebakan}

\begin{variasi}
Titik potong sumbu-$x$ (butuh menyelesaikan $3x^2-x+2=0$; cek $D=1-24<0$ $\Rightarrow$
tidak memotong sumbu-$x$).
\end{variasi}

% ---------------- SOAL 6.1.10 ----------------
\soalno{6.1.10}{waktu tinggi maksimum (proyektil)}
\begin{soalbox}{}
Sebuah batu dilempar, tingginya $h(t)=20t-5t^2$. Kapan mencapai tinggi
maksimum?
\end{soalbox}

\jawab
\langkah{1} Tulis dalam bentuk baku: $h(t)=-5t^2+20t$, jadi $a=-5,\ b=20$.
Karena $a<0$, ada maksimum.
\langkah{2} Waktu puncak: $t=-\dfrac{b}{2a}=-\dfrac{20}{2(-5)}=\dfrac{20}{10}=2$.
\jadi tinggi maksimum dicapai pada $t=2$ detik. (Tinggi maksnya
$h(2)=40-20=20$ m.)

\begin{jebakan}
\begin{itemize}[leftmargin=*,topsep=2pt]
  \item \textbf{Salah menyusun $a,b$.} Susun $-5t^2+20t$ (koefisien $t^2$ adalah
  $-5$).
  \item \textbf{Mengira ``kapan menyentuh tanah'' = puncak.} Menyentuh tanah saat
  $h=0$ ($t=0$ atau $t=4$); puncak di tengahnya, $t=2$.
\end{itemize}
\end{jebakan}

\begin{variasi}
Menanyakan tinggi maksimum (substitusi $t=2$), atau kapan menyentuh tanah
($h=0$), atau berapa lama di atas $15$ m (pertidaksamaan).
\end{variasi}

% ############################################################
\section*{Subbab 6.2 — Persamaan, Pertidaksamaan, dan Penerapan}
\addcontentsline{toc}{section}{Subbab 6.2 — Persamaan \& Pertidaksamaan}
% ############################################################

\begin{ingat}{Rumus $abc$, Vieta, dan pertidaksamaan}
Akar: $x_{1,2}=\dfrac{-b\pm\sqrt{b^2-4ac}}{2a}$. Vieta:
$x_1+x_2=-\dfrac ba$, $x_1x_2=\dfrac ca$. \\
Pertidaksamaan kuadrat: cari akar, gambar parabola, lalu baca daerahnya --- untuk
$a>0$, nilai \emph{positif} di \emph{luar} akar, \emph{negatif} di \emph{antara}
akar.
\end{ingat}

% ---------------- SOAL 6.2.1 ----------------
\soalno{6.2.1}{rumus abc}
\begin{soalbox}{}
Selesaikan $x^2-7x+10=0$ dengan rumus $abc$.
\end{soalbox}

\jawab
\langkah{1} $a=1,\ b=-7,\ c=10$. $D=(-7)^2-4(1)(10)=49-40=9$.
\langkah{2} $x=\dfrac{-b\pm\sqrt D}{2a}=\dfrac{7\pm\sqrt9}{2}=\dfrac{7\pm3}{2}$.
\langkah{3} $x=\dfrac{10}{2}=5$ atau $x=\dfrac{4}{2}=2$.
\jadi $x=5$ atau $x=2$.

\begin{jebakan}
\begin{itemize}[leftmargin=*,topsep=2pt]
  \item \textbf{Salah tanda $-b$.} $b=-7$, jadi $-b=+7$.
\end{itemize}
\end{jebakan}

\begin{variasi}
Bandingkan dengan pemfaktoran $(x-2)(x-5)$; rumus $abc$ berlaku untuk yang sulit
difaktorkan.
\end{variasi}

% ---------------- SOAL 6.2.2 ----------------
\soalno{6.2.2}{Vieta dasar}
\begin{soalbox}{}
Jika $x_1,x_2$ akar $x^2-3x+1=0$, hitung $x_1+x_2$ dan $x_1x_2$.
\end{soalbox}

\jawab
\langkah{1} $a=1,\ b=-3,\ c=1$.
\langkah{2} Jumlah akar: $x_1+x_2=-\dfrac ba=-\dfrac{-3}{1}=3$.
\langkah{3} Hasil kali: $x_1x_2=\dfrac ca=\dfrac11=1$.
\jadi $x_1+x_2=3$ dan $x_1x_2=1$.

\begin{jebakan}
\begin{itemize}[leftmargin=*,topsep=2pt]
  \item \textbf{Lupa tanda minus pada jumlah.} $x_1+x_2=-\frac ba$; di sini
  $-(-3)=3$.
\end{itemize}
\end{jebakan}

\begin{variasi}
Lanjutkan menghitung $x_1^2+x_2^2=(x_1+x_2)^2-2x_1x_2=9-2=7$ (kombinasi simetris).
\end{variasi}

% ---------------- SOAL 6.2.3 ----------------
\soalno{6.2.3}{ekspresi simetris akar}
\begin{soalbox}{}
Hitung $\dfrac{1}{x_1}+\dfrac{1}{x_2}$ untuk $2x^2-5x+3=0$.
\end{soalbox}

\begin{ingat}{Ubah ke jumlah \& hasil kali}
$\dfrac{1}{x_1}+\dfrac{1}{x_2}=\dfrac{x_1+x_2}{x_1x_2}$. Ekspresi simetris apa pun
bisa ditulis lewat $x_1+x_2$ dan $x_1x_2$ (Vieta), tanpa mencari akar.
\end{ingat}

\jawab
\langkah{1} Vieta: $x_1+x_2=-\dfrac{-5}{2}=\dfrac52$; $x_1x_2=\dfrac32$.
\langkah{2} $\dfrac{1}{x_1}+\dfrac{1}{x_2}=\dfrac{x_1+x_2}{x_1x_2}
=\dfrac{5/2}{3/2}=\dfrac53$.
\jadi $\dfrac{1}{x_1}+\dfrac{1}{x_2}=\dfrac53$.

\begin{jebakan}
\begin{itemize}[leftmargin=*,topsep=2pt]
  \item \textbf{Menjumlahkan kebalikan secara terpisah} setelah mencari akar ---
  boros. Vieta jauh lebih cepat.
  \item \textbf{Lupa $a=2$} sehingga $x_1+x_2=5$ (salah); bagi dengan $a$.
\end{itemize}
\end{jebakan}

\begin{variasi}
$x_1^2+x_2^2$, $\frac{x_1}{x_2}+\frac{x_2}{x_1}$, $(x_1-x_2)^2=(x_1+x_2)^2-4x_1x_2$.
\end{variasi}

% ---------------- SOAL 6.2.4 ----------------
\soalno{6.2.4}{pertidaksamaan (di antara akar)}
\begin{soalbox}{}
Tentukan HP dari $x^2-4\le0$.
\end{soalbox}

\jawab
\langkah{1} Akar $x^2-4=0\Rightarrow(x-2)(x+2)=0\Rightarrow x=\pm2$.
\langkah{2} $a=1>0$ (buka atas), nilai $\le0$ (di bawah/pada sumbu) berada
\emph{di antara} akar.
\jadi HP $=\{x\mid -2\le x\le2\}$.

\begin{center}
\begin{tikzpicture}[scale=0.8]
\draw[->](-3.5,0)--(3.5,0) node[right]{$x$};
\foreach \x in {-2,2}{\fill (\x,0) circle(2.5pt);}
\draw[very thick,merah](-2,0)--(2,0);
\node[below] at (-2,-0.15){$-2$};\node[below] at (2,-0.15){$2$};
\node[merah,above] at (0,0.1){HP: $-2\le x\le2$};
\end{tikzpicture}
\end{center}

\begin{jebakan}
\begin{itemize}[leftmargin=*,topsep=2pt]
  \item \textbf{Menulis $x\le2$ saja.} $x^2\le4$ berarti $|x|\le2$, jadi
  \emph{dua sisi}: $-2\le x\le2$.
\end{itemize}
\end{jebakan}

\begin{variasi}
$x^2-4\ge0$ membalik jawaban jadi $x\le-2$ atau $x\ge2$ (di luar akar).
\end{variasi}

% ---------------- SOAL 6.2.5 ----------------
\soalno{6.2.5}{pertidaksamaan (di luar akar)}
\begin{soalbox}{}
Tentukan HP dari $x^2+x-12\ge0$.
\end{soalbox}

\jawab
\langkah{1} Akar: $x^2+x-12=0\Rightarrow(x+4)(x-3)=0\Rightarrow x=-4,\ 3$.
\langkah{2} $a=1>0$ (buka atas), nilai $\ge0$ berada \emph{di luar} akar
(termasuk titik akarnya).
\jadi HP $=\{x\mid x\le-4\ \text{atau}\ x\ge3\}$.

\begin{jebakan}
\begin{itemize}[leftmargin=*,topsep=2pt]
  \item \textbf{Menggabung jadi satu pertidaksamaan} ($-4\le x\le3$) --- itu
  untuk ``$\le0$''. Untuk ``$\ge0$'' daerahnya \emph{terpisah}.
\end{itemize}
\end{jebakan}

\begin{variasi}
Bentuk dengan $a<0$ (mis.\ $-x^2+x+12\ge0$) membalik logika ``luar/dalam''.
\end{variasi}

% ---------------- SOAL 6.2.6 ----------------
\soalno{6.2.6}{menyusun PK dari akar}
\begin{soalbox}{}
Susun persamaan kuadrat berakar $2$ dan $-5$.
\end{soalbox}

\begin{ingat}{PK dari jumlah \& hasil kali}
PK dapat ditulis $x^2-(x_1+x_2)x+(x_1x_2)=0$.
\end{ingat}

\jawab
\langkah{1} Jumlah akar: $2+(-5)=-3$. Hasil kali: $2\cdot(-5)=-10$.
\langkah{2} Susun: $x^2-(-3)x+(-10)=0\Rightarrow x^2+3x-10=0$.
\jadi $x^2+3x-10=0$ (cek: $(x-2)(x+5)=x^2+3x-10$ \checkmark).

\begin{jebakan}
\begin{itemize}[leftmargin=*,topsep=2pt]
  \item \textbf{Salah tanda koefisien $x$.} Rumus memakai $-(x_1+x_2)$; di sini
  $-(-3)=+3$.
\end{itemize}
\end{jebakan}

\begin{variasi}
Akar irasional/pecahan; atau PK yang akarnya $2\times$ akar PK lain (transformasi
akar).
\end{variasi}

% ---------------- SOAL 6.2.7 ----------------
\soalno{6.2.7}{satu akar diketahui}
\begin{soalbox}{}
Jika satu akar $x^2+px+12=0$ adalah $3$, tentukan $p$ dan akar lainnya.
\end{soalbox}

\jawab
\textbf{Cara A (substitusi).} Karena $3$ akar, ia memenuhi persamaan:
$3^2+p(3)+12=0\Rightarrow 9+3p+12=0\Rightarrow 3p=-21\Rightarrow p=-7$.

\textbf{Akar lain (Vieta).} Hasil kali akar $=\dfrac{c}{a}=12$. Jika satu akar
$3$: $3\cdot x_2=12\Rightarrow x_2=4$.
\jadi $p=-7$ dan akar lainnya $4$. (Cek jumlah: $3+4=7=-p$ \checkmark.)

\begin{jebakan}
\begin{itemize}[leftmargin=*,topsep=2pt]
  \item \textbf{Mencari akar lain lewat rumus $abc$} padahal Vieta ($x_1x_2=c$)
  jauh lebih cepat.
\end{itemize}
\end{jebakan}

\begin{variasi}
Diberi hasil/jumlah akar tertentu untuk menentukan parameter; akar berkebalikan.
\end{variasi}

% ---------------- SOAL 6.2.8 ----------------
\soalno{6.2.8}{akar kembar cari parameter}
\begin{soalbox}{}
Tentukan $m$ agar $x^2-mx+9=0$ berakar kembar.
\end{soalbox}

\jawab
\langkah{1} Syarat kembar: $D=0$. Di sini $a=1,\ b=-m,\ c=9$:
\[
D=(-m)^2-4(1)(9)=m^2-36=0.
\]
\langkah{2} $m^2=36\Rightarrow m=\pm6$.
\jadi $m=6$ atau $m=-6$.

\begin{jebakan}
\begin{itemize}[leftmargin=*,topsep=2pt]
  \item \textbf{Hanya menulis $m=6$.} $m^2=36$ punya \emph{dua} solusi
  ($\pm6$).
  \item \textbf{$(-m)^2=-m^2$?} Salah; $(-m)^2=m^2$ (kuadrat selalu positif).
\end{itemize}
\end{jebakan}

\begin{variasi}
Syarat $D>0$ (dua akar berbeda) memberi $m<-6$ atau $m>6$; $D<0$ memberi
$-6<m<6$.
\end{variasi}

% ---------------- SOAL 6.2.9 ----------------
\soalno{6.2.9}{pertidaksamaan buka bawah}
\begin{soalbox}{}
Selesaikan $-x^2+3x>0$.
\end{soalbox}

\jawab
\langkah{1} Kalikan $-1$ (\textbf{balik tanda}): $x^2-3x<0$.
\langkah{2} Faktorkan: $x(x-3)<0\Rightarrow$ akar $0$ dan $3$.
\langkah{3} $a=1>0$, nilai $<0$ berada \emph{di antara} akar.
\jadi HP $=\{x\mid 0<x<3\}$.

\begin{tips}
Alternatif tanpa membalik: $-x^2+3x=x(3-x)>0$. Hasil kali positif bila kedua
faktor sepertanda: $0<x<3$. Sama hasilnya.
\end{tips}

\begin{jebakan}
\begin{itemize}[leftmargin=*,topsep=2pt]
  \item \textbf{Lupa membalik tanda} saat mengalikan $-1$. $>$ berubah jadi $<$.
\end{itemize}
\end{jebakan}

\begin{variasi}
$-x^2+3x\le0$ (di luar akar: $x\le0$ atau $x\ge3$).
\end{variasi}

% ---------------- SOAL 6.2.10 ----------------
\soalno{6.2.10}{akar berbanding (Vieta HOTS)}
\begin{soalbox}{}
Akar $x^2-6x+k=0$ berbanding $1:2$. Tentukan $k$.
\end{soalbox}

\begin{ingat}{Misalkan akar sesuai perbandingan}
Kalau akar berbanding $1:2$, tulis akarnya $t$ dan $2t$. Lalu pakai Vieta:
jumlah $=-\frac ba$, hasil kali $=\frac ca$.
\end{ingat}

\jawab
\langkah{1} Misalkan akarnya $t$ dan $2t$.
\langkah{2} Jumlah akar $=-\dfrac{-6}{1}=6$: $t+2t=3t=6\Rightarrow t=2$.
\langkah{3} Maka akarnya $2$ dan $4$. Hasil kali $=k$: $k=2\cdot4=8$.
\jadi $k=8$. (Cek: $x^2-6x+8=(x-2)(x-4)$, akar $2:4=1:2$ \checkmark.)

\begin{jebakan}
\begin{itemize}[leftmargin=*,topsep=2pt]
  \item \textbf{Menganggap akarnya $1$ dan $2$.} Perbandingan $1:2$ berarti
  $t:2t$, bukan nilai persisnya $1$ dan $2$.
\end{itemize}
\end{jebakan}

\begin{variasi}
Akar berbanding $2:3$, atau berselisih tertentu ($x_2-x_1=4$), atau satu akar
$k\times$ akar lain.
\end{variasi}

% ============================================================
\begin{tips}
\textbf{Ringkasan strategi Bab 6.}
\begin{enumerate}[leftmargin=*,topsep=2pt]
  \item \emph{Grafik dulu:} tanda $a$ (arah buka), sumbu simetri
  $x=-\frac{b}{2a}$, puncak, titik potong. Gambar mengungkap banyak jawaban.
  \item \emph{Diskriminan} $D=b^2-4ac$ menentukan banyak akar; $D=0$ untuk akar
  kembar (dipakai mencari parameter).
  \item \emph{Vieta} ($x_1+x_2=-\frac ba,\ x_1x_2=\frac ca$) menyelesaikan
  ekspresi simetris tanpa mencari akar.
  \item \emph{Pertidaksamaan:} cari akar, lihat arah buka --- $a>0$: positif di
  luar akar, negatif di antara. \textbf{Balik tanda} bila mengali negatif.
\end{enumerate}
\end{tips}

\end{document}
